% Mixture of Experts (MoE) models have emerged as the de-facto architecture for scaling up language models without significantly increasing the computational cost.\wentao{xinyu: It is not easy for people to directly understand this sentence. You can see a tradeoff (more fine-grained: better performance; less fine-grained: better efficiency). Maybe FlashMoBA writes something similar. I would also prefer something like system efficiency vs model performance} 
% Existing MoE methods optimize system efficiency or model architecture independently. We show that as MoE models get more granular and sparser, they become more memory-bound, and jointly optimizing the model architecture and kernel design leads to a significant improvement training throughput for MoE models. 

Mixture of Experts (MoE) models have emerged as the de facto architecture for scaling up language models without significantly increasing the computational cost. Recent MoE models demonstrate a clear trend towards high expert granularity (smaller expert intermediate dimension) and higher sparsity (constant number of activated experts with a higher number of total experts), which improve model quality per FLOP. However, fine-grained MoEs suffer from increased activation memory footprint and reduced hardware efficiency due to higher IO costs, while sparser MoEs suffer from wasted computations due to padding in Grouped GEMM kernels. In response, we propose a memory-efficient algorithm to compute the forward and backward passes of MoEs with minimal activation caching for the backward pass. We also design GPU kernels that overlap memory IO with computation, benefiting all MoE architectures. Finally, we propose a novel ``token rounding'' method that minimizes the wasted compute due to padding in Grouped GEMM kernels. As a result, our method \Algname\ \textbf{reduces activation memory by 45\% and achieves a 1.86x compute throughput improvement} on Hopper GPUs compared to ScatterMoE's BF16 MoE kernel for a fine-grained 7B MoE. Concretely, \Algname\ on 64 H100s achieves a training throughput of 213 billion tokens per day, comparable to ScatterMoE's 225 billion tokens per day on 96 H100s for a 7B MoE model training with FSDP-2 using the lm-engine codebase\footnote{\small \url{https://github.com/open-lm-engine/lm-engine}}. \textbf{On Blackwell GPUs, \Algname~also achieves a 25\% and 15\% relative speedup on the forward and backward pass respectively compared to a highly optimized DeepGEMM baseline on OLMoE-sized 7B MoE models.} Under high MoE sparsity settings, our tile-aware token rounding algorithm yields an additional \textbf{1.16x} speedup on kernel execution time compared to vanilla top-$K$ routing while maintaining similar downstream performance on Hopper GPUs. We open-source all our kernels\footnote{\small \url{https://github.com/Dao-AILab/sonic-moe}} to enable faster MoE model training.

% Mixture of Experts (MoE) models have emerged as the de-facto architecture for scaling up language models without significantly increasing the computational cost. Existing MoE methods optimize system efficiency or model architecture independently. We show that as MoE models get more granular and sparser, they become more memory-bound, and jointly optimizing the algorithms and the kernel design leads to a major improvement in MoE training throughput. We first propose a memory-efficient algorithm to compute the forward and backward of MoE with minimal activation saved. We then design GPU kernels that overlap memory IO latency with compute, benefiting all MoE architectures. Finally, we propose a novel "token rounding" method that minimizes the wasted compute brought by tile quantization. As a result, our method SNaX reduces 45% activation memory and has 1.87x compute throughput improvement on Hopper GPUs compared to state-of-the-art BF16 MoE kernel for a fine-grained 7B MoE. Concretely, SNaX on 64 H100s achieves almost the same total throughput as ScatterMoE on 96 H100s for a 7B MoE model training with token-choice rounding while training with FSDP-2. Under high MoE sparsity settings, our tile-aware token rounding algorithm yields an additional 1.18x speedup on kernel execution time compared to vanilla top-K routing while maintaining similar downstream performance.
